\documentclass[a4paper, 12pt]{article}

\usepackage[portuges]{babel}
\usepackage[utf8]{inputenc}
\usepackage{amsmath}
\usepackage{indentfirst}
\usepackage{graphicx}
\usepackage{multicol,lipsum}
\usepackage{xurl}
\usepackage[colorlinks=true, urlcolor=blue]{hyperref}

\begin{document}
%\maketitle

\begin{titlepage}
	\begin{center}
	
	%\begin{figure}[!ht]
	%\centering
	%\includegraphics[width=2cm]{c:/ufba.jpg}
	%\end{figure}

		\Large{UNESP - IBILCE}\\
		\large{Departamento de Ciência da Computação e Estatísticas}\\ 
		\large{Disciplina de Compiladores - Bacharelado em Ciência da Computação}\\ 
		\vspace{15pt}
        \vspace{95pt}
        \textbf{\LARGE{Analisador Léxico - Trabalho 1}}\\
		%\title{{\large{Título}}}
		\vspace{3,5cm}
	\end{center}
	
	\begin{flushleft}
		\begin{tabbing}
			Alunos: Gabriel Martins Brum e João Carlos Gonçalves Nobre \\
			Professora: Profa. Dra. Renata Spolon Lobato
	\end{tabbing}
 \end{flushleft}
	\vspace{1cm}
	
	\begin{center}
		\vspace{\fill}
			 Outubro\\
		 2026
			\end{center}
\end{titlepage}
%%%%%%%%%%%%%%%%%%%%%%%%%%%%%%%%%%%%%%%%%%%%%%%%%%%%%%%%%%%



\newpage
\pagenumbering{arabic}
% % % % % % % % % % % % % % % % % % % % % % % % % % %

\section{Descrição da Linguagem Regular}

A linguagem regular desenvolvida para este analisador léxico é baseada nos requisitos do Trabalho 1 e composta pelas seguintes definições:

\begin{itemize}
    \item \textbf{Identificadores:} Iniciam obrigatoriamente com uma letra ou sublinhado (\verb|[a-zA-Z_]|), podendo ser seguidos por letras, números ou sublinhados adicionais.
    \item \textbf{Números Inteiros:} Formados por uma sequência contínua de um ou mais dígitos (\verb|[0-9]+|).
    \item \textbf{Números Reais:} Definidos por uma parte inteira de um ou mais dígitos, seguida de um ponto separador (\verb|.|) e uma parte fracionária (\verb|[0-9]+|).
    \item \textbf{Operadores:} A linguagem reconhece os operadores aritméticos de soma (\verb|+|), subtração (\verb|-|), multiplicação (\verb|*|) e divisão (\verb|/|), além do operador de atribuição (\verb|=|).
    \item \textbf{Pontuação e Delimitadores:} São reconhecidos os símbolos de vírgula (\verb|,|), ponto (\verb|.|), ponto e vírgula (\verb|;|), dois pontos (\verb|:|), além de chaves (\verb|{|, \verb|}|) e parênteses (\verb|(|, \verb|)|) para aninhamento de instruções.
    \item \textbf{Comentários:} Delimitados estritamente por colchetes (\verb|[| e \verb|]|). O analisador reconhece e ignora formalmente qualquer conteúdo interno aos mesmos, sem gerar tokens.
    \item \textbf{Tratamento de Erros e Espaços:} Espaços em branco, quebras de linha e tabulações são ignorados, enquanto quebras de linha incrementam o contador de rastreamento. Qualquer caractere não previsto pelas regras descritas acima é capturado e gera o aviso de \textit{Erro léxico}, indicando o caractere inválido e a linha de ocorrência.
\end{itemize}

\section{Manual do Usuário}

Para facilitar a compilação e a validação do analisador léxico, o projeto foi estruturado com scripts Bash de automação. É necessário possuir o compilador \textit{GCC} e a ferramenta \textit{Flex} devidamente instalados no sistema.

\subsection*{Compilação Rápida}
Para gerar o analisador léxico a partir do código fonte, basta utilizar o script de compilação fornecido no diretório:
\begin{verbatim}
$ chmod +x compile.sh
$ ./compile.sh
\end{verbatim}
O script \texttt{compile.sh} encarrega-se automaticamente de invocar o Flex sobre o arquivo \texttt{lexer.l} e, na sequência, utilizar o GCC para converter o código intermediário \texttt{lex.yy.c} no executável final do analisador.

\subsection*{Execução e Rotina de Testes}
O analisador pode ser utilizado isoladamente apontando para qualquer arquivo de texto, sendo o arquivo \texttt{input1.txt} o explicitamente solicitado no enunciado do trabalho, que não contém erros léxicos. A execução é realizada da seguinte forma:
\begin{verbatim}
$ ./lexer input1.txt
\end{verbatim}

Para validar todos os limites e tratamentos de erros discutidos durante o desenvolvimento, criamos um ambiente de teste em lote. Basta utilizar o script \texttt{run\_tests.sh}:
\begin{verbatim}
$ chmod +x run_tests.sh
$ ./run_tests.sh
\end{verbatim}
Este script percorrerá automaticamente os casos de teste rigorosos (de \texttt{input1.txt} até \texttt{input5.txt}), imprimindo as respectivas quebras de linha e saídas no terminal. 

Para validar a integridade da implementação, consulte o arquivo auxiliar \texttt{expected\_output.txt}. Ele contém o mapeamento e o gabarito detalhado de todos os tokens e erros léxicos que devem ser apresentados no terminal durante a execução da bateria de testes.

\subsection*{Repositório \textit{Github}}

Caso seja de interesse, todo o código fonte e arquivos para a realização deste trabalho estão disponível em \url{https://github.com/gabrielmbrum/graduacao/tree/main/8th-semester/compiladores/trabalho1/}.

\end{document}
